% =====================================================================
%  Plan Trip Module — Study & Defense Guide  (STANDALONE, not part of the thesis)
%  Compile with pdfLaTeX on Overleaf. Explanations in Vietnamese; the
%  practice answers are in English so they can be rehearsed for the jury.
% =====================================================================
\documentclass[12pt,a4paper]{article}
\usepackage[utf8]{inputenc}
\usepackage{lmodern}
\usepackage[T5]{fontenc}
\usepackage[vietnamese]{babel}
\usepackage{amsmath}
\usepackage[margin=2.4cm]{geometry}
\usepackage{enumitem}
\usepackage{xcolor}
\usepackage[hidelinks]{hyperref}

\setlength{\parskip}{0.5em}
\setlength{\parindent}{0pt}
\newcommand{\en}[1]{\textit{(#1)}}            % English term gloss
\newcommand{\Q}[1]{\medskip\noindent\textbf{Q:} \textit{#1}\par}
\newcommand{\A}{\noindent\textbf{A:} }

\title{\textbf{HanoiGO --- Module Plan Trip}\\[4pt]
\large Tài liệu giải thích \& luyện bảo vệ trước hội đồng}
\author{Nguyen Duy Hung --- 22BA13154}
\date{}

\begin{document}
\maketitle
\hrule
\medskip

\textbf{Cách dùng tài liệu:} phần 1--8 giải thích module bằng tiếng Việt (kèm
thuật ngữ tiếng Anh trong ngoặc để bạn nhớ từ vựng). Phần 9 là bộ câu hỏi luyện
tập: câu hỏi và câu trả lời mẫu viết bằng tiếng Anh để bạn nói trực tiếp trước
hội đồng. Các con số đều là kết quả thật từ benchmark, không bịa.

\tableofcontents
\medskip\hrule

% =====================================================================
\section{Module Plan Trip làm gì?}

Module này nhận một danh sách địa điểm và vài thiết lập, rồi tự sinh ra một
\textbf{lịch trình nhiều ngày} \en{multi-day itinerary}.

\begin{itemize}[nosep]
  \item \textbf{Đầu vào} \en{input}: các địa điểm đã chọn; số ngày đi; giờ bắt
  đầu / kết thúc mỗi ngày; khung nghỉ trưa; điểm xuất phát (khách sạn).
  \item \textbf{Đầu ra} \en{output}: với mỗi ngày, một danh sách điểm \textbf{theo
  thứ tự}, kèm giờ đến \en{arrival} và giờ đi \en{departure}. Điểm nào không kịp
  giờ mở cửa thì báo là \textit{unscheduled}.
\end{itemize}

Ví von dễ hiểu: module đóng vai một hướng dẫn viên, xếp lịch sao cho khách
\textbf{đi được nhiều điểm nhất}, \textbf{ít di chuyển nhất}, và \textbf{không
tới nơi lúc đã đóng cửa}.

% =====================================================================
\section{Bài toán cốt lõi: TSPTW}

Đây là bài toán \textbf{Người giao hàng với khung giờ} \en{Travelling Salesman
Problem with Time Windows --- TSPTW}.

\begin{itemize}[nosep]
  \item \textbf{TSP} (bài toán người giao hàng): tìm thứ tự thăm các điểm sao cho
  tổng quãng đường / thời gian di chuyển nhỏ nhất. Đây là bài toán
  \textbf{NP-hard} --- số thứ tự khả dĩ tăng theo giai thừa $n!$.
  \item \textbf{Time Windows} (khung giờ): mỗi điểm chỉ mở trong một khoảng giờ
  nhất định (ví dụ Lăng Bác 07:30--10:30). Phải thăm \textbf{trong} khung đó.
\end{itemize}

Vì sao khó: phải đồng thời (1) chọn thứ tự tốt, (2) tôn trọng giờ mở/đóng của
từng điểm, (3) chừa giờ nghỉ trưa, và (4) không vượt quá giờ kết thúc ngày. Một
thứ tự ``ngắn đường'' nhưng tới Lăng Bác sau 10:30 là vô dụng.

% =====================================================================
\section{Quy trình 6 bước (pipeline) --- chi tiết từng bước}

Hãy hình dung cả module như một \textbf{trợ lý sắp lịch} nhận vào danh sách
địa điểm bạn chọn và trả ra lịch trình từng ngày. Trợ lý này không làm mọi thứ
trong một bước, mà đi qua một chuỗi giai đoạn, giai đoạn sau dọn dẹp và tinh
chỉnh kết quả của giai đoạn trước. Dưới đây là toàn bộ hành trình của dữ liệu,
từ lúc bạn bấm ``Tạo lịch trình'' tới lúc nhận về kết quả.

\textbf{Giai đoạn chuẩn bị --- lấy thông tin từng điểm.} Trước tiên hệ thống
tra trong cơ sở dữ liệu để biết: điểm này ở toạ độ nào, mở cửa những ngày nào
trong tuần, mở--đóng lúc mấy giờ, và nên dành bao nhiêu phút để tham quan. Đây
chỉ là bước ``lấy hồ sơ'' cho từng địa điểm, chưa tính toán gì cả.

\textbf{Bước 1 --- Loại bỏ những điểm vô vọng ngay từ đầu.} Với mỗi địa điểm,
hệ thống hỏi một câu đơn giản: \textit{``Trong toàn bộ những ngày của chuyến đi
này, có ngày nào điểm này mở cửa không?''}. Nếu câu trả lời là \textbf{không}
--- ví dụ bạn đi đúng 2 ngày cuối tuần mà một bảo tàng chỉ mở các ngày trong
tuần --- thì điểm đó bị loại ngay từ vòng gửi xe, không tốn công tính toán gì
thêm cho nó. Đây gọi là nhóm \textbf{``không khả thi''}, khác với nhóm
\textbf{``không xếp được''} sẽ xuất hiện ở các bước sau (điểm có thể mở cửa
đúng ngày, nhưng lịch trình không còn chỗ hoặc không kịp giờ).

\textbf{Bước 2 --- Chia điểm thành từng ngày theo vị trí địa lý.} Đây là bước
trả lời câu hỏi \textit{``ngày nào nên đi những điểm nào?''}, hoàn toàn dựa vào
\textbf{khoảng cách trên bản đồ}, chưa quan tâm giờ giấc. Cách làm giống như
việc bạn nhìn bản đồ và khoanh vùng: những điểm nằm gần nhau được gom vào cùng
một khoanh, số khoanh bằng đúng số ngày của chuyến đi. Kỹ thuật dùng ở đây là
\textbf{K-Means++} --- một thuật toán gom cụm rất phổ biến, với cải tiến quan
trọng là chọn các ``tâm cụm'' ban đầu \textbf{dàn trải khắp khu vực} thay vì
chọn ngẫu nhiên hoàn toàn, để tránh việc 2 tâm cụm vô tình rơi sát nhau ngay từ
đầu. Việc chọn tâm ban đầu vẫn có yếu tố ngẫu nhiên, nhưng dùng một công thức
sinh số \textbf{luôn cho cùng một kết quả với cùng một tập điểm đầu vào} ---
điều này quan trọng để bạn chạy đi chạy lại benchmark mà không lo kết quả nhảy
lung tung mỗi lần.

Sau khi chia xong, có một bước \textbf{cân bằng lại}: mỗi ngày chỉ được chứa
tối đa 5 điểm (giới hạn hợp lý cho một ngày tham quan thực tế). Nếu một cụm bị
dồn quá 5 điểm, hệ thống sẽ đẩy \textbf{điểm lạc lõng nhất} của cụm đó (điểm xa
tâm cụm nhất) sang cụm khác, nhưng chỉ đẩy sang cụm nào \textbf{gần điểm đó
nhất} và không làm cụm đích trở nên quá rời rạc. Nói cách khác: nếu bắt buộc
phải bớt người trong một nhóm đông, ta đẩy người đang đứng xa nhất sang nhóm
gần họ nhất, chứ không đẩy bừa.

\textbf{Bước 3 --- Sửa lỗi ``đúng ngày đóng cửa''.} Sau bước gom cụm theo vị
trí, có thể xảy ra tình huống trớ trêu: một điểm bị xếp vào đúng ngày mà nó
đóng cửa (ví dụ nhiều bảo tàng ở Hà Nội nghỉ thứ Hai, mà điểm đó lại rơi vào
cụm ứng với thứ Hai). Bước này quét lại toàn bộ các ngày, phát hiện những
trường hợp như vậy, rồi tìm một ngày khác trong chuyến mà điểm đó \textbf{mở
cửa và còn chỗ trống} để chuyển sang.

\textbf{Bước 4 --- Hỏi Goong Maps về thời gian di chuyển thật.} Đến đây, hệ
thống đã biết ngày nào đi những điểm nào, nhưng chưa biết \textbf{thứ tự} thăm
trong ngày. Muốn xếp thứ tự cho đúng, cần biết đi từ điểm này sang điểm kia mất
bao lâu \textbf{ngoài đời thật} (có hồ chắn đường, phố một chiều, tắc đường...)
chứ không thể chỉ đo đường chim bay. Vì vậy hệ thống gọi một lần duy nhất tới
Goong Maps, hỏi cùng lúc: \textit{``từ điểm xuất phát của người dùng tới từng
điểm mất bao lâu?''} và \textit{``giữa mọi cặp điểm với nhau mất bao lâu?''}.
Gọi một lần cho cả chuyến (thay vì gọi riêng cho từng ngày) giúp tiết kiệm thời
gian chờ và tránh bị giới hạn tần suất gọi API. Nếu Goong không phản hồi được
(lỗi mạng, hết hạn mức), hệ thống không bị treo mà tự động chuyển sang
\textbf{ước lượng bằng công thức đường chim bay} kết hợp tốc độ trung bình xe
máy 30\,km/h.

\textbf{Bước 5 --- Xếp thứ tự thăm trong từng ngày.} Đây là bước quan trọng
nhất về mặt thuật toán, được trình bày riêng và chi tiết ở Mục~4 (Trái tim
thuật toán). Tóm tắt: với ngày có từ 8 điểm trở xuống, hệ thống \textbf{thử tất
cả các cách sắp xếp có thể} để tìm ra cách tốt nhất một cách chắc chắn; với
ngày nhiều điểm hơn (hiếm khi xảy ra vì mỗi ngày đã giới hạn 5 điểm), hệ thống
dùng một cách \textbf{ước lượng nhanh} thay thế.

\textbf{Bước 6 --- Cập nhật lại giờ giấc theo GPS thật.} Thứ tự vừa chọn ở bước
5 được tính với một giả định đơn giản hoá: \textit{``coi như đi tới điểm đầu
tiên trong ngày mất 0 phút''}. Đến bước này, hệ thống thay con số 0 phút giả
định đó bằng \textbf{thời gian di chuyển thật} từ chỗ người dùng đang đứng
(khách sạn, hoặc vị trí họ chọn) tới điểm đầu tiên, rồi \textbf{đẩy lùi toàn bộ
các mốc giờ phía sau} tương ứng --- giống như khi chuyến bay đầu ngày bị hoãn
thì cả lịch trình cả ngày bị lùi theo. Nếu việc lùi giờ này khiến một điểm nào
đó rơi vào sau giờ đóng cửa, điểm đó buộc phải bị loại khỏi ngày, và mọi mốc
giờ được tính lại từ đầu cho những điểm còn ở lại.

\textbf{Bước 7 --- Cố tìm chỗ cho những điểm đã bị rớt.} Những điểm bị loại ở
các bước trước (đóng cửa đúng ngày ban đầu, không kịp giờ sau khi lùi lịch...)
chưa bị bỏ cuộc hẳn. Hệ thống thử: với mỗi điểm bị rớt, xét từng ngày mà điểm
đó \textbf{còn mở cửa}, rồi thử nhét nó vào từng \textbf{khe hở} giữa hai điểm
liên tiếp (hoặc đầu/cuối ngày) --- hễ tìm được một khe hở mà nhét vào không làm
lệch giờ của các điểm xung quanh thì chèn ngay tại đó. Chỉ khi không còn khe hở
nào phù hợp ở bất kỳ ngày nào, điểm đó mới chính thức bị đánh dấu \textbf{không
xếp được} và báo cho người dùng biết.

\textbf{Bước 8 --- Xếp thứ tự lại một lần cuối.} Sau bước 7, danh sách điểm của
mỗi ngày có thể đã thay đổi so với lúc bước 5 chạy lần đầu (mất bớt điểm rớt,
có thêm điểm được chèn vào). Vì vậy thứ tự tính ở bước 5 chưa chắc còn là thứ
tự tốt nhất cho \textbf{danh sách điểm cuối cùng} của ngày đó. Bước này chạy
lại đúng cách xếp thứ tự ở bước 5 (thử hết mọi khả năng, hoặc ước lượng nhanh
tuỳ số điểm) một lần nữa cho danh sách đã chốt, để đảm bảo thứ tự cuối cùng
trả về cho người dùng là tốt nhất có thể. Có một \textbf{cơ chế phòng hộ}: nếu
việc xếp lại này chẳng may làm rớt mất một điểm vốn dĩ đã có mặt (vì giờ GPS
thật đẩy lịch trễ hơn dự kiến), hệ thống sẽ \textbf{giữ nguyên bản trước đó}
thay vì chấp nhận bản mới --- nghĩa là bước xếp lại cuối cùng này chỉ có thể
giữ nguyên hoặc làm tốt hơn, không bao giờ khiến kết quả tệ đi.

\textbf{Ý chính cần nhớ khi trình bày:} các bước đầu (chuẩn bị dữ liệu, loại bỏ
điểm vô vọng, gom cụm theo vị trí, lấy thời gian di chuyển thật) chỉ là
\textbf{heuristic nhanh} để dọn dẹp bài toán cho gọn; hai lần xếp thứ tự (bước
5 và bước 8) mới là nơi \textbf{thuật toán chính xác} thực sự phát huy tác
dụng; các bước ở giữa (cập nhật GPS thật, chèn lại điểm rớt) là phần
\textbf{đưa lịch trình từ lý thuyết về sát với thực tế}.

% =====================================================================
\section{Trái tim thuật toán: Tham lam vs Tìm kiếm chính xác}

\subsection{Trước tiên: ``chi phí'' được định nghĩa như thế nào?}

Cả hai cách xếp thứ tự dưới đây (tham lam và tìm kiếm chính xác) đều cần một
cách để so sánh ``phương án này tốt hơn phương án kia bao nhiêu''. Thước đo đó
gọi là \textbf{hàm chi phí} \en{cost function}, và nó \textbf{không phải một
công thức có sẵn đi kèm thuật toán} --- đây là phần do chính người xây dựng
module này tự định nghĩa, để phù hợp với bài toán du lịch (chứ một thuật toán
tham lam hay tìm kiếm chính xác ``nguyên bản'' trong sách giáo khoa chỉ nói
\textit{``chọn phương án chi phí thấp nhất''}, không quy định chi phí đó tính
ra sao).

Công thức được chọn ở đây, tính theo đơn vị giây, là:
\[
\text{Chi phí} \;=\; \big(\text{thời gian di chuyển tính bằng giây}\big) \times 2
\;+\; \big(\text{thời gian đứng chờ tính bằng giây}\big) \times 1
\]
Vì thời gian chờ trong hệ thống được lưu theo \textbf{phút} còn thời gian di
chuyển theo \textbf{giây}, nên khi viết bằng biến của code, phút chờ phải nhân
với 60 để quy về cùng đơn vị giây trước khi cộng.

\textbf{Vì sao lại nhân thời gian di chuyển với 2?} Đây là một \textbf{lựa
chọn heuristic có chủ đích}, không phải một định lý hay công thức chuẩn: 1
giây ngồi trên xe di chuyển được coi là ``tốn'' gấp đôi 1 giây đứng chờ tại chỗ
(chờ trước giờ mở cửa). Lý do trực giác: di chuyển gắn với hao tốn nhiên liệu,
công sức, và rủi ro (tắc đường, thời tiết), trong khi chờ chỉ đơn thuần là mất
thời gian. Nếu hội đồng hỏi ``con số 2 này có căn cứ khoa học nào không'',
câu trả lời trung thực là: \textbf{đây là trọng số chọn theo kinh nghiệm/thử
nghiệm} (không phải suy ra từ một mô hình lý thuyết), và việc tinh chỉnh trọng
số này một cách có hệ thống (ví dụ bằng thử nghiệm A/B hoặc tối ưu tham số) là
một hướng cải tiến trong tương lai --- không nên bịa ra một cơ sở lý thuyết
không tồn tại.

\subsection{Tham lam (Greedy Nearest Neighbor --- GNN)}

Tham lam nghĩa là: tại mỗi thời điểm, chỉ nhìn về phía trước \textbf{đúng một
bước}, và chọn ngay điểm nào có chi phí thấp nhất theo công thức ở trên, mà
không hề cân nhắc xem lựa chọn đó ảnh hưởng ra sao tới các bước xa hơn về sau.
Đây là một chiến lược rất phổ biến trong khoa học máy tính vì nó nhanh và dễ
cài đặt, nhưng nhược điểm cố hữu của nó là \textbf{thiển cận} --- một lựa chọn
có vẻ tốt ngay lúc này có thể đẩy cả kế hoạch vào ngõ cụt sau đó vài bước.

\subsection{Cái bẫy Ba Đình --- ví dụ kinh điển cho thấy tham lam thất bại}

Hãy tưởng tượng bạn xuất phát từ một khách sạn cạnh Hồ Gươm lúc 8 giờ sáng, dự
định thăm 5 điểm trong khu Ba Đình. Trong 5 điểm đó, Lăng Bác đóng cửa lúc
10:30 và một bảo tàng gần đó đóng cửa lúc 12:00; ba điểm còn lại mở cửa cả
ngày.

Một người lập kế hoạch theo kiểu tham lam sẽ nghĩ rất đơn giản: \textit{``ngay
bây giờ, điểm nào gần khách sạn nhất?''} --- và chọn ngay Bảo tàng Phụ nữ vì nó
ở sát khách sạn. Sau đó lại tiếp tục hỏi câu tương tự cho bước kế tiếp, vô tình
la cà qua các điểm mở cửa cả ngày trước, để rồi khi cuối cùng nhớ ra phải ghé
Lăng Bác thì đồng hồ đã chỉ quá 10:30 --- \textbf{Lăng Bác đã đóng cửa, bị bỏ
lỡ hoàn toàn}. Kết quả: chỉ thăm được 3 trên 5 điểm, dù quãng đường đi mỗi bước
đều là ``ngắn nhất tại thời điểm đó''.

Cách làm đúng lại ngược trực giác ban đầu: phải \textbf{ưu tiên ghé điểm nào
đóng cửa sớm nhất trước}, dù nó không phải điểm gần nhất. Đi theo thứ tự Bảo
tàng Hồ Chí Minh rồi tới Lăng Bác trước, sau đó mới vòng qua các điểm còn lại,
sẽ thăm được trọn vẹn cả 5 điểm. Đây chính là bài học cốt lõi: \textbf{``gần
nhất ngay bây giờ'' không đồng nghĩa với ``tốt nhất cho cả ngày''}.

\subsection{Tìm kiếm chính xác (Exact / Brute-force search)}

Thay vì đoán từng bước một như tham lam, cách tiếp cận này \textbf{thử tất cả
mọi thứ tự có thể có}, tính chi phí cho từng thứ tự, rồi mới chọn ra thứ tự tốt
nhất. Với $n$ điểm cần thăm, số thứ tự khả dĩ là $n!$ (n giai thừa) --- ví dụ 5
điểm có $5! = 120$ cách sắp xếp khác nhau, 8 điểm có tới $8! = 40{,}320$ cách.

Việc chọn ra thứ tự ``tốt nhất'' trong số đó dựa trên \textbf{hai tiêu chí, xét
lần lượt theo thứ tự ưu tiên}:
\begin{enumerate}[nosep]
  \item \textbf{Tiêu chí 1 (quan trọng nhất): thứ tự nào thăm được nhiều điểm
  nhất.} Một thứ tự xếp đủ 5/5 điểm luôn thắng một thứ tự chỉ xếp được 4/5,
  \textbf{bất kể} thứ tự 4 điểm kia có chi phí thấp hơn bao nhiêu.
  \item \textbf{Tiêu chí 2 (chỉ xét khi hai thứ tự hoà nhau ở tiêu chí 1):
  thứ tự nào có tổng chi phí (di chuyển + chờ, theo công thức ở trên) thấp
  hơn.}
\end{enumerate}
Sở dĩ phải có hai tầng thay vì chỉ so chi phí, là vì nếu chỉ xét chi phí đơn
thuần, một thứ tự có thể ``lách luật'' bằng cách \textbf{bỏ hẳn một điểm khó
xếp} (điểm có khung giờ hẹp) để chi phí trông có vẻ thấp hơn --- nhưng thực tế
người dùng bị mất một điểm tham quan. Ưu tiên ``thăm đủ điểm'' lên trên ``chi
phí thấp'' đảm bảo hệ thống không đánh đổi trải nghiệm người dùng chỉ để có
con số đẹp.

\textbf{Ví dụ số minh hoạ hai tầng ưu tiên này hoạt động ra sao:} giả sử ngày
hôm đó có 3 điểm ứng viên, xuất phát từ một điểm ở ngoại thành lúc 8:00 sáng.

\begin{center}
\renewcommand{\arraystretch}{1.15}
\begin{tabular}{|l|c|c|c|}
\hline
\textbf{Điểm} & \textbf{Cách xuất phát} & \textbf{Giờ mở cửa} & \textbf{Thời gian tham quan} \\ \hline
A --- xa nhất & 15\,km & mở cả ngày (08:00--18:00) & 60 phút \\ \hline
B --- gần nhất & 3\,km & chỉ mở từ 10:00 & 45 phút \\ \hline
C --- ở giữa & 8\,km & mở cả ngày (08:00--18:00) & 30 phút \\ \hline
\end{tabular}
\end{center}

Nhìn thoáng qua, ai cũng sẽ nghĩ nên đi B trước vì B gần nhất. Nhưng thử tính
hai phương án:

\textit{Phương án ``đi gần trước'' (B $\to$ A $\to$ C):} 8:00 xuất phát, chỉ mất
khoảng 6 phút để tới B (3\,km, tốc độ trung bình 30\,km/h) --- nhưng B tận
10:00 mới mở cửa, nên phải \textbf{đứng chờ khoảng 114 phút}. Riêng khoản chờ
này đã cộng vào chi phí một lượng rất lớn (114 phút $\times$ 60 giây/phút = 6\,840
đơn vị chi phí), chưa kể thời gian di chuyển sau đó.

\textit{Phương án ``đi xa trước, hợp giờ hơn'' (A $\to$ C $\to$ B):} 8:00 xuất
phát, đi thẳng tới A (15\,km, mất khoảng 30 phút), tham quan xong ghé C, rồi
cuối cùng mới tới B. Vì lúc này đã hơn 10:00, B \textbf{đã mở cửa sẵn, gần như
không phải chờ}. Tổng thời gian di chuyển của phương án này cao hơn phương án
trên (vì đi xa trước), nhưng phần chi phí do phải chờ gần như bằng 0.

So sánh hai phương án: phương án ``đi gần trước'' tuy tổng quãng đường di
chuyển ngắn hơn, nhưng chi phí chờ đợi (114 phút $\times$ 60) áp đảo hoàn toàn
phần chênh lệch quãng đường (chỉ khoảng vài trăm giây). Vì vậy \textbf{tìm kiếm
chính xác sẽ chọn phương án A $\to$ C $\to$ B} --- tức là điểm xa nhất (A) được
xếp \textbf{đầu tiên}, còn điểm gần nhất (B) lại bị đẩy xuống \textbf{cuối
cùng}. Đây chính xác là kiểu hiện tượng ``điểm gần lại không được xếp đầu tiên''
mà một người quan sát bằng mắt thường trên bản đồ dễ cảm thấy vô lý --- nhưng
thực chất là kết quả tính toán đúng đắn của thuật toán, không phải lỗi.

\textbf{Vì sao chấp nhận được việc thử tất cả $n!$ cách, dù đây là con số tăng
rất nhanh?} Vì mỗi ngày trong ứng dụng này luôn bị giới hạn ở quy mô rất nhỏ:
tối đa 5 điểm theo thiết kế, và một trần an toàn tuyệt đối là 8 điểm (vượt quá
sẽ chuyển sang dùng cách tham lam để tránh bùng nổ số phép thử). Ngay tại
trần 8 điểm, tổng số cách sắp xếp là $8! = 40{,}320$ --- máy tính hiện đại thử
hết toàn bộ số này chỉ trong \textbf{chưa tới 1 mili-giây}. Trong khi đó, một
lần gọi Goong Maps để lấy thời gian di chuyển thật đã mất từ 200 đến 500 mili-giây.
Nói cách khác, ở quy mô bài toán của module này, việc thử ``tất cả khả năng''
không hề tốn kém --- nó rẻ hơn hàng trăm lần so với việc gọi một API bên ngoài.

\subsection{Một lần sửa lỗi quan trọng: tính đến chặng di chuyển đầu tiên}

Ban đầu, bước tìm kiếm chính xác giả định một điều đơn giản hoá: \textit{``đi
từ điểm xuất phát tới điểm đầu tiên trong ngày mất 0 giây''}. Giả định này chỉ
đúng trên lý thuyết. Khi bước sau đó (cập nhật GPS thật) cộng thêm thời gian
di chuyển thật từ khách sạn (chẳng hạn khoảng 14 phút), toàn bộ các mốc giờ
trong ngày bị lùi lại theo, và một điểm có khung giờ hẹp gần giữa trưa có thể
bất ngờ bị đẩy quá giờ đóng cửa rồi bị loại --- dù lúc tìm kiếm chính xác chọn
thứ tự, thứ tự đó trông vẫn hoàn toàn hợp lệ.

Cách khắc phục: cho bước tìm kiếm chính xác \textbf{biết trước} thời gian di
chuyển thật từ điểm xuất phát tới từng điểm ứng viên (thông tin này đã có sẵn
từ bước gọi Goong Maps trước đó, không cần gọi thêm), để nó chọn ra một thứ tự
\textbf{đã tính đến độ trễ thực tế ngay từ đầu}, thay vì chọn một thứ tự chỉ
đẹp trên lý thuyết rồi vỡ trận khi áp thời gian thật vào. Sau khi sửa, ví dụ
benchmark ở mục dưới thăm đủ trọn vẹn 5/5 điểm.

% =====================================================================
\section{Vì sao chọn từng kỹ thuật (design rationale)}

\begin{itemize}[nosep]
  \item \textbf{K-Means++ để chia ngày, không phải để xếp thứ tự.} Chia điểm
  theo cụm địa lý đảm bảo mỗi ngày gọn trong một khu, đỡ chạy lòng vòng giữa các
  quận xa nhau.
  \item \textbf{Haversine cho bước gom cụm} \en{straight-line distance}: gom cụm
  chỉ cần khoảng cách tương đối, nên dùng công thức đường chim bay tính
  \textbf{offline, miễn phí, tức thì}.
  \item \textbf{Goong cho bước xếp thứ tự}: xếp lịch theo giờ cần
  \textbf{thời gian đi đường thật} (tránh hồ, phố một chiều), nên dùng Goong
  Distance Matrix.
  \item \textbf{Dự phòng Haversine + retry}: nếu Goong lỗi/giới hạn tần suất, thử
  lại 3 lần (exponential backoff) rồi quay về ước lượng Haversine --- hệ không
  bao giờ chết.
  \item \textbf{Đệm tiếp cận 10 phút \& nghỉ trưa}: mỗi điểm cộng 10 phút đệm cho
  việc tiếp cận thực tế --- từ chỗ xuống xe/taxi tới cổng, mua hoặc quét vé, định
  vị --- trước khi bắt đầu thăm. Đây là con số bảo thủ, \textbf{không phụ thuộc
  loại phương tiện} (đi xe máy, Grab hay đi bộ đều cần), giúp giờ in ra trung
  thực. Nếu giờ thăm đụng khung trưa thì đẩy sang sau trưa.
\end{itemize}

% =====================================================================
\section{Đánh giá: chứng minh tốt hơn cách thủ công}

So sánh lịch của module \textbf{HanoiGO} với hai cách du khách hay tự làm:
\textbf{By-District} (gom theo quận, thăm theo thứ tự liệt kê) và \textbf{Random}
(ngẫu nhiên, trung bình 200 lần). Cả ba dùng \textbf{cùng dữ liệu và cùng bộ
kiểm tra giờ giấc}, chỉ khác chiến lược.

\begin{center}
\renewcommand{\arraystretch}{1.2}
\begin{tabular}{|l|l|c|c|c|c|}
\hline
\textbf{Kịch bản} & \textbf{Phương pháp} & \textbf{Điểm thăm} & \textbf{Di chuyển} & \textbf{Chờ} & \textbf{Độ gọn} \\
 & & & \textbf{(phút)} & \textbf{(phút)} & \textbf{(km)} \\ \hline
Case 1 (1 ngày,    & HanoiGO     & \textbf{5/5}   & 25  & 0   & 0.80 \\ \cline{2-6}
5 điểm,            & By-District & 4/5            & 19  & 0   & 0.91 \\ \cline{2-6}
bẫy Ba Đình)       & Random      & 3.6/5          & 25  & 0   & 0.96 \\ \hline
Case 2 (3 ngày,    & HanoiGO     & \textbf{12/12} & \textbf{100} & 0 & \textbf{0.35} \\ \cline{2-6}
12 điểm,           & By-District & 12/12          & 112 & 0   & 1.53 \\ \cline{2-6}
nhiều quận)        & Random      & 10.2/12        & 136 & 0   & 1.93 \\ \hline
Case 3 (2 ngày,    & HanoiGO     & \textbf{8/8}   & 33  & 0   & 0.61 \\ \cline{2-6}
8 điểm, đóng cửa   & By-District & 5/8            & 31  & 0   & 0.36 \\ \cline{2-6}
thứ Hai)           & Random      & 6.3/8          & 44  & 0   & 0.77 \\ \hline
Case 4 (1 ngày,    & HanoiGO     & \textbf{5/5}   & 27  & \textbf{0}   & 0.37 \\ \cline{2-6}
5 điểm, mở         & By-District & 3/5            & 15  & 404 & 0.20 \\ \cline{2-6}
buổi chiều)        & Random      & 4.3/5          & 20  & 257 & 0.32 \\ \hline
\end{tabular}
\end{center}

\textbf{Đọc kết quả:} các con số trên dùng \textbf{thời gian đi đường thật của
Goong} (ảnh chụp đông cứng để tái lập được). HanoiGO \textbf{thăm đủ mọi điểm
người dùng chọn ở cả 4 case}. Ở Case 2, lợi thế chính là \textbf{độ gọn}: ngày
của nó gọn gấp khoảng 4 lần (0.35 so với 1.53 km), đồng thời tổng di chuyển vẫn
thấp nhất (100 so với 112 và 136 phút) --- dù khi dùng đường thật, khoảng cách
với By-District thu hẹp lại (đường thật phạt các ngày trải rộng về địa lý, nhưng
By-District ở Case 2 vẫn giữ được 12/12 nên travel không chênh nhiều). Case 3 và
Case 4 mới là nơi module bỏ xa cách thủ công, vì chúng kiểm tra hai ràng buộc
giờ giấc mà cách thủ công bỏ qua (xem mục dưới).

\subsection*{Case 1 --- Bẫy Ba Đình (1 ngày, 5 điểm)}

\textbf{Bối cảnh:} chuyến 1 ngày trong khu Ba Đình. Lăng Bác đóng cửa
lúc 10:30 và một bảo tàng gần đó đóng lúc 12:00; ba điểm còn lại mở cả ngày.
Xuất phát từ khách sạn gần Hồ Gươm lúc 8:00.

\begin{itemize}[nosep]
  \item \textbf{HanoiGO} nhận ra Lăng Bác đóng sớm nhất, xếp nó \textbf{đầu ngày}
  rồi ghé bảo tàng trước 12:00, cuối cùng thăm đủ \textbf{5/5} điểm.
  \item \textbf{By-District} chọn ngay điểm gần khách sạn nhất (Bảo tàng Phụ nữ),
  lòng vòng các điểm mở-cả-ngày rồi \textbf{tới Lăng Bác sau 10:30 --- rớt}
  (4/5). Random rớt trung bình $\approx 1.4$ điểm (3.6/5).
\end{itemize}

\textbf{Lưu ý thiết kế quan trọng (dễ bị hỏi):} phiên bản đầu của module dùng
\textbf{tham lam (GNN)} cho mọi ngày và cũng bỏ lỡ Lăng Bác theo cách tương tự.
Benchmark này chính là thứ phát hiện ra lỗi đó --- sau đó tôi đổi sang
\textbf{tìm kiếm chính xác} cho các ngày nhỏ ($\le 8$ điểm). Đây là ví dụ điển
hình về ``benchmark dẫn đến cải tiến thiết kế'', không chỉ là bảng số đẹp.

\textbf{Câu trả lời nếu bị hỏi về di chuyển cao hơn By-District (25 vs.\ 19
phút):} HanoiGO thăm 5 điểm trong khi By-District chỉ thăm 4, nên đương nhiên
tổng di chuyển dài hơn một chặng. Đây là đánh đổi hợp lý: một điểm bổ sung = 6
phút di chuyển thêm.

\subsection*{Case 2 --- Nhiều quận, nhiều ngày (3 ngày, 12 điểm)}

\textbf{Bối cảnh:} chuyến 3 ngày với 12 điểm trải rộng qua nhiều quận Hà Nội,
không có điểm nào đóng đặc biệt sớm và chuyến đi không rơi vào thứ Hai. Đây là
kịch bản ``không có bẫy giờ giấc'' --- thử thách đến từ việc \textbf{chia ngày
hợp lý về địa lý}.

\begin{itemize}[nosep]
  \item \textbf{HanoiGO} và \textbf{By-District} đều thăm đủ \textbf{12/12}
  (Random rớt $\approx 2$ điểm: 10.2/12).
  \item Lợi thế của HanoiGO nằm ở \textbf{độ gọn}: 0.35 km so với 1.53 km của
  By-District --- gọn hơn gấp $\approx 4$ lần. Các quận hành chính không khớp
  với mạng lưới đường thực tế, nên gom theo quận tạo ra những ngày trải rộng
  địa lý; K-Means++ gom theo khoảng cách thực nên mỗi ngày gọn về mặt bản đồ.
  \item Tổng di chuyển cũng thấp nhất (100 phút so với 112 và 136), dù khoảng
  cách với By-District thu hẹp khi dùng đường thật thay vì đường chim bay (vì
  đường thật phạt các ngày trải rộng địa lý, nhưng By-District ở case này vẫn
  giữ được 12/12 nên travel không chênh nhiều).
\end{itemize}

\textbf{Câu trả lời nếu bị hỏi ``By-District cũng 12/12, vậy HanoiGO hơn gì?'':}
Thăm đủ điểm mới chỉ là điều kiện cần; \textbf{trải nghiệm hành trình} cũng quan
trọng. Ngày gọn hơn 4 lần nghĩa là khách đỡ phải di chuyển lòng vòng giữa các
khu xa nhau trong một ngày. Lợi thế thực sự của Case 2 là \textbf{chất lượng
hành trình mỗi ngày}, không phải số điểm đến. Độ chênh travel time nhỏ (100 vs
112) vì đường thật co khoảng cách lại; nhưng độ chênh độ gọn (0.35 vs 1.53) vẫn
rất lớn và duy trì ý nghĩa trong thực tế.

\subsection*{Case 3 --- Đóng cửa theo thứ trong tuần}

\textbf{Bối cảnh:} chuyến 2 ngày, \textbf{ngày 1 rơi vào thứ Hai}. Trong 8 điểm
chọn, có 3 điểm đóng cửa thứ Hai (Bảo tàng Hồ Chí Minh, Hoàng thành Thăng Long,
Bảo tàng Phụ nữ --- nhiều bảo tàng ở Hà Nội nghỉ đầu tuần là chuyện có thật).

\begin{itemize}[nosep]
  \item \textbf{HanoiGO} đọc danh sách ngày-mở \en{openDays} của từng điểm, nên
  \textbf{đẩy cả 3 điểm đóng sang ngày 2} (thứ Ba) và thăm đủ \textbf{8/8}.
  \item \textbf{By-District} không quan tâm hôm đó là thứ mấy, để 3 điểm đó ở
  ngày thứ Hai rồi \textbf{rớt cả 3} (5/8). Random rớt trung bình $\approx 2$
  điểm (6.3/8).
\end{itemize}

\textbf{Lưu ý trung thực (rất dễ bị hỏi):} ở bảng, By-District ``gọn'' hơn
(0.36 so với 0.61) --- nhưng \textbf{chỉ vì nó đã rớt 3 điểm xa}, nên độ gọn được
đo trên ít điểm hơn. HanoiGO chịu kém gọn một chút để \textbf{giữ đủ điểm}. Đừng
để con số độ gọn đánh lừa: tiêu chí quan trọng nhất ở đây là số điểm thăm được.

\textbf{Ý nghĩa bảo vệ:} đây chính là lý do app bắt người dùng \textbf{nhập ngày
đi} --- để biết thứ trong tuần và không xếp khách tới nơi đúng hôm đóng cửa.

\subsection*{Case 4 --- Điểm chỉ mở buổi chiều}

\textbf{Bối cảnh:} 1 ngày, có \textbf{Nhà hát Múa rối Thăng Long mở từ 15:00}
(suất chiều/tối) cùng 4 điểm mở buổi sáng.

\begin{itemize}[nosep]
  \item \textbf{HanoiGO} xếp nhà hát \textbf{cuối ngày}, lấp buổi sáng bằng 4 điểm
  kia, nên \textbf{không phải chờ} (0 phút) và thăm đủ 5/5.
  \item \textbf{By-District} giữ nhà hát ở vị trí cố định (đầu danh sách), tới nơi
  lúc sáng rồi \textbf{ngồi chờ tới 15:00 = 404 phút}; thời gian mất vào việc chờ
  khiến nó \textbf{rớt 2 điểm} (3/5). Random chờ 257 phút.
\end{itemize}

\textbf{Vì sao HanoiGO tự biết làm vậy:} hàm chi phí phạt cả thời gian chờ
($\text{Cost} = t_{\text{travel}}\times 2 + t_{\text{wait}}$), nên một điểm mở
muộn tự bị đẩy về cuối. Case này cũng là lý do bảng có \textbf{cột ``Chờ''}: khi
mọi điểm đều mở cả ngày thì cột này bằng 0 và vô nghĩa, nhưng khi có điểm mở muộn
thì nó cho thấy rõ giá trị của việc xếp theo giờ mở cửa.

% =====================================================================
\section{Hạn chế \& đánh đổi (phải nói trung thực)}

\begin{itemize}[nosep]
  \item \textbf{Không tối ưu toàn cục} \en{not globally optimal}: tìm kiếm chính
  xác chỉ tối ưu \textbf{trong từng ngày} với tập điểm đã cho; còn việc
  \textbf{chia điểm vào ngày} (K-Means) là heuristic, nên tổng thể chưa chắc tối
  ưu tuyệt đối.
  \item \textbf{Có thể tách hai điểm gần nhau sang hai ngày}: khi hai điểm sát
  nhau cùng đóng cửa buổi trưa (ví dụ một điểm đóng 11:00, một điểm đóng 13:00),
  mỗi điểm cần trọn một buổi sáng và khung nghỉ trưa chen vào giữa, nên không thể
  xếp chung một ngày. Hệ \textbf{ưu tiên tính khả thi về giờ giấc hơn độ gọn địa
  lý} --- đây là đánh đổi của cách gom cụm, không phải lỗi.
  \item \textbf{Trần 8 điểm/ngày}: trên 8 điểm phải quay về GNN để tránh bùng nổ
  giai thừa (thực tế hiếm vì cap 5).
  \item \textbf{Chỉ tính theo xe máy} \en{motorbike mode}; chưa hỗ trợ phương
  tiện công cộng.
  \item \textbf{Phụ thuộc Goong}; số liệu travel time là ảnh chụp tại một thời
  điểm, sẽ đổi theo giao thông.
  \item \textbf{Benchmark dùng fixture cố định}, không phải toàn bộ cơ sở dữ
  liệu, nên mô tả chuyến đi tiêu biểu chứ không phải mọi trường hợp.
\end{itemize}

\subsection*{Ví dụ thực tế: chuyến \texttt{test123} (2 ngày, 7 điểm)}

Một chuyến người dùng tạo thật minh hoạ rõ hạn chế ``tách hai điểm gần nhau sang
hai ngày'':

\begin{itemize}[nosep]
  \item \textbf{Ngày 1}: Lăng Bác \en{mở 07:30--11:00} $\rightarrow$ Cột cờ Hà
  Nội $\rightarrow$ Vườn thú Thủ Lệ.
  \item \textbf{Ngày 2}: Bảo tàng B52 \en{mở 09:00--13:00} $\rightarrow$ Bảo tàng
  Mỹ thuật $\rightarrow$ Phố tàu hoả $\rightarrow$ Hồ Gươm.
\end{itemize}

\textbf{Nghịch lý:} Cột cờ Hà Nội (Ngày 1) và Bảo tàng Mỹ thuật (Ngày 2) chỉ cách
nhau \textbf{0,35 km} nhưng lại bị xếp vào hai ngày khác nhau; tương tự Cột cờ và
Phố tàu hoả cách \textbf{0,49 km} cũng bị tách. Trong khi đó các điểm \emph{trong
cùng một ngày} lại xa nhau hơn (tới hơn 2 km).

\textbf{Vì sao:} hai ``điểm neo'' đóng cửa buổi trưa --- Lăng Bác (đóng 11:00) và
B52 (đóng 13:00) --- \textbf{không thể chung một ngày}. Mỗi điểm cần trọn một buổi
sáng, mà khung nghỉ trưa 11:00--13:00 chen vào giữa khiến không kịp thăm điểm thứ
hai trước giờ đóng. Vì thế chúng buộc phải neo ở hai ngày khác nhau. Sau đó
K-Means gắn các điểm ``linh hoạt'' (mở cả ngày) vào cụm của điểm neo \emph{gần
nhất}, nên Cột cờ dính theo cụm Lăng Bác còn Mỹ thuật dính theo cụm B52 --- dù hai
điểm này sát nhau.

\textbf{Kết luận:} lịch \textbf{vẫn hợp lệ} --- không vi phạm giờ giấc, không bỏ
lỡ điểm nào --- nhưng \textbf{chưa gọn nhất về địa lý}. Đây đúng là đánh đổi
``khả thi theo giờ $>$ gọn theo địa lý'', không phải lỗi. Hướng cải tiến: sau khi
cố định các điểm neo theo giờ, chạy thêm một bước \textbf{hoán đổi cục bộ}
\en{local search} để dồn các điểm linh hoạt sát nhau về cùng một ngày.

% =====================================================================
\section{Bảng thuật ngữ nhanh}

\begin{description}[leftmargin=0pt,style=nextline,itemsep=2pt]
  \item[TSPTW] Travelling Salesman Problem with Time Windows --- xếp thứ tự thăm
  có ràng buộc khung giờ.
  \item[K-Means++] thuật toán gom cụm điểm theo vị trí; ``++'' là cách chọn tâm
  cụm ban đầu cho trải đều hơn.
  \item[Heuristic] cách giải nhanh, gần đúng, không đảm bảo tối ưu.
  \item[Exact / Brute-force] thử mọi khả năng để chắc chắn tìm ra lời giải tốt
  nhất (trong phạm vi nhỏ).
  \item[Time window] khung giờ mở--đóng của một địa điểm.
  \item[Haversine] công thức tính khoảng cách đường chim bay trên mặt cầu.
  \item[GPS cascade] cộng thời gian đi thật từ điểm xuất phát rồi dồn dịch các
  mốc giờ phía sau.
\end{description}

% =====================================================================
\section{Câu hỏi luyện tập bảo vệ (English)}

Câu trả lời mẫu viết bằng tiếng Anh để luyện nói. Hãy tập trả lời ngắn gọn,
tự tin, và \textbf{thừa nhận hạn chế khi cần}.

\Q{Brute-force search is factorial, $O(n!)$. Why is it acceptable in your planner?}
\A Each day is capped at a very small number of places --- at most five by design,
eight as a hard limit. With at most $8! = 40{,}320$ orderings, one day is solved
in under a millisecond, far cheaper than a single Goong API call (200--500 ms).
The factorial cost only matters for large $n$, which never happens within one day.

\Q{Why not use a dedicated TSP solver such as Google OR-Tools, or a metaheuristic
like genetic algorithms or simulated annealing?}
\A Those tools are designed for problems with tens or hundreds of nodes, where
exact search is impossible. My sub-problem has at most eight nodes, so the exact
search already returns the guaranteed best order instantly, with no extra
dependency and fully deterministic results. Using a heavier solver would add
complexity without improving the answer.

\Q{Why split the problem into clustering then ordering, instead of solving it as
one big optimisation?}
\A Solving multi-day TSPTW jointly is very expensive. Splitting it keeps each
part tractable: K-Means++ quickly assigns places to days by geography, and the
exact search then optimises each small day. It is a divide-and-conquer trade-off:
slightly less than a global optimum, but fast and good enough for a web app.

\Q{Why use Haversine for clustering but the Goong API for sequencing? Isn't that
inconsistent?}
\A They serve different needs. Clustering only needs relative closeness, so the
free, offline Haversine distance is enough and keeps it fast. Sequencing schedules
against real opening hours, so it needs accurate road travel times, which is why
it uses Goong. Using the right tool for each step is a deliberate design choice.

\Q{What exactly is a ``time window'', and how does the planner respect it?}
\A A time window is the open--close interval of a place, for example the
Mausoleum from 07:30 to 10:30. For every candidate order, the planner simulates
arrival, waiting, and departure times; a stop is kept only if the visit finishes
before both its closing time and the daily end time, otherwise it is dropped.

\Q{How do you handle the lunch break?}
\A The lunch break (default 11:00--13:00) is a blocked window. If a visit would
overlap it, the arrival is shifted to after lunch and the following stops are
recomputed. It is user-configurable.

\Q{Why do you add a fixed 10-minute buffer at every stop?}
\A It is a realistic access buffer, not a parking fee. At each site a visitor
still needs time to get from the drop-off point to the entrance, buy or scan a
ticket, and orient themselves before the visit actually starts. Ten minutes is a
conservative constant that keeps the printed arrival and departure times honest.
It is independent of the transport mode, so it holds whether the tourist rides a
motorbike, takes a Grab, or walks --- which matters because most foreign visitors
do not self-drive.

\Q{What happens if the Goong API fails or is rate-limited?}
\A The system retries up to three times with exponential backoff. If it still
fails, it falls back to estimating travel time from the Haversine distance and an
average motorbike speed of 30 km/h, so the planner always returns a result.

\Q{How do you prove the planner is actually better than planning by hand?}
\A I ran a benchmark comparing HanoiGO against two manual strategies, By-District
and Random, on the same places, the same settings, and the same feasibility
check, with frozen real Goong road travel times for reproducibility. Across all
four scenarios HanoiGO scheduled every place the user selected. In the multi-day
case its days were about four times more compact (0.35 versus 1.53 km) with the
lowest total travel time, and in the two constrained cases it respected Monday
closures and an afternoon-only opening time that the manual methods ignored.

\Q{How did you keep the benchmark fair between methods?}
\A All three methods received identical inputs and used the exact same
time-window feasibility function, so the only variable was the planning strategy.
The HanoiGO arm called the real production service, and the Random method was
averaged over 200 seeds to avoid cherry-picking a lucky case.

\Q{In Case 2, By-District also visited all 12 places. So what is the real
advantage of your planner?}
\A Visiting all places is only half the goal; the experience matters too. For the
same 12 places, my planner kept far more compact days, about four times tighter
(0.35 versus 1.53 km), and still the lowest total travel time (100 against 112
minutes). Administrative districts do not match the real road network, so grouping
by geography is what keeps each day tight; the travel margin is smaller once real
road durations are used, but the compactness gap stays large.

\Q{Did you find any flaw during evaluation, and how did you fix it?}
\A Yes. The benchmark revealed that the greedy heuristic chose each day's places
short-sightedly and stranded the early-closing Mausoleum. I changed small days to
use the exact search, and also made it account for the real first-leg travel time
from the start point, after which the planner scheduled every place.

\Q{What is the GPS cascade and why is it needed?}
\A The ordering step initially assumes the day starts exactly at the first stop.
The GPS cascade injects the real travel time from the user's start point to the
first stop and shifts all later times accordingly, then drops or re-checks any
stop that no longer fits. It makes the printed schedule match reality.

\Q{Is your solution globally optimal?}
\A No, and I do not claim that. The ordering within each day is exact and optimal
for that day's set of places, but assigning places to days uses K-Means, which is
a heuristic. So the result is high quality and feasible, but not guaranteed
globally optimal --- a reasonable trade-off for an interactive web app.

\Q{Two places in the itinerary are very close together, yet the planner put them
on different days. Is that a bug?}
\A No. Day assignment starts from geography via K-Means, but it is overridden by
time-window feasibility. When two nearby places both close around midday --- say
one at 11:00 and another at 13:00 --- each needs a full morning slot before it
shuts, and the fixed lunch break sits between them, so they cannot share the same
morning. The planner therefore anchors them on separate days even though they are
physically close. It deliberately prioritises feasibility over pure geographic
compactness; this is the clustering trade-off listed in my limitations. To avoid
confusing a viewer during a live demo, I pick example trips whose early-closing
sites do not collide in this way.

\Q{What if a user selects 50 places for a 2-day trip?}
\A Clustering would put about 25 places per day, far above the 5-place cap, so
capacity rebalancing moves the overflow and some places become unscheduled and
are reported. Days above eight places fall back to the greedy heuristic. In
practice the UI encourages a realistic number of places per day.

\Q{How are times represented in the code?}
\A Clock times are stored as integer minutes from midnight, which makes the
window arithmetic simple, while travel durations from Goong are in seconds. This
avoids time-zone and date-object pitfalls.

\Q{What would you improve in future work?}
\A Public-transport and walking modes besides motorbike; a joint multi-day
optimisation instead of clustering-then-ordering; and expanding the dataset
beyond Hanoi. I would also re-run the benchmark on a larger, randomly sampled set
of real trips to strengthen the statistical claim.

\Q{Your planner asks for a travel date. Why does the exact date matter?}
\A Because opening hours depend on the day of the week; several Hanoi museums
close on Mondays. The planner reads each place's open days and moves a place that
is closed on a chosen day to another day where it is open. In the Monday test
scenario it scheduled all eight places, while grouping by district left three on
the closed day and dropped them.

\Q{In one of your cases, By-District shows a better compactness number than
HanoiGO. Doesn't that mean it planned better there?}
\A No. By-District looks more compact only because it dropped three far-apart
places, so its compactness is measured over fewer points. HanoiGO kept all eight,
including the distant ones, so its days are slightly less tight. Scheduling every
place the user chose matters more than a compactness value computed over a smaller
set, so I report visited count as the primary metric.

\Q{How does the planner deal with a place that only opens in the afternoon?}
\A The cost function penalises waiting time, so a late-opening place is naturally
pushed toward the end of the day. In the afternoon-only test, the planner filled
the morning with the other places and visited the 15:00 theatre last with zero
waiting, whereas a fixed order that kept the theatre early idled for over six
hours and lost two places to the wasted time.

\end{document}
